\begin{abstract}
接吻数是离散几何中的经典问题，目前仅在少数维度得到精确解。
近年来，人工智能通过大规模数值与组合搜索发现了更大的接吻构型，
但如何将搜索结果发展为可推广的数学构造与严格证明，仍是一项挑战。
本研究利用自主多智能体科研系统Qiushi Engine，取得十九个维数的新下界，
覆盖25、27、32—39、43、45及49—55维。
这些构造来自不同的结构机制，包括接触层的协同运动、带标签的方向复用、
联合支持交换、局部符号码替换、跨壳格构造、低重叠等距像以及球面设计矩证书。
所得结果包括
$K(25)\ge197580$、$K(27)\ge201567$、$K(38)\ge591900$、
$K(43)\ge2553792$、$K(45)\ge7380090$及$K(55)\ge53301140$。
Qiushi Engine自主完成构造搜索、数学分析与计算验证，
各项最终结果均落实为明确的数学论证和可独立检查的有限证书。
这些成果展示了自主科研系统如何超越固定表述下的优化，
在大规模探索中发现新的数学表示与构造。

\par\smallskip\noindent\textbf{关键词}\quad
自主数学发现；接吻数；几何变形；常重码；格提升；符号设计；球面设计
\end{abstract}
